% ======================================================================
% Chapitre : Théorie de l'Apprentissage et Minimisation du Risque Empirique
% ======================================================================

Dans les chapitres précédents, nous avons présenté différentes méthodes d'apprentissage basées sur la minimisation du risque empirique (ERM, \emph{Empirical Risk Minimization}). Or la seule analyse de la convergence de ces méthodes que nous ayons faite est la décomposition biais-variance du risque ainsi que la décomposition de son erreur en deux : l'erreur \emph{d'approximation} et l'erreur \emph{d'estimation}. 

Dans la seconde partie de ce chapitre, nous nous attardons justement sur l'erreur faite sur le risque lors de l'apprentissage. Pour cela, nous développons une nouvelle théorie basée sur la complexité de l'espace de recherche $ \mathcal{F}$. 

Mais avant, nous développons le concept de \emph{convexification du risque} dans la première partie. En effet, pour certaines fonction de pertes (notamment la \emph{0-1 loss}) nous ne disposons pas de "bonnes" propriétés analytiques pour utiliser des algorithmes d'optimisation. Nous présentons ainsi une méthode permettant d'y remédier. 

% ======================================================================
% Convexification du risque et substituts convexes

\section{Convexification du risque et substituts convexes}

% -------------------------------------------------
\subsection{Reparamétrisation et identité de marge}

En classification binaire, on cherche à prédire une étiquette parmi deux classes $\mathcal{Y} = \{-1, 1\}$ via une fonction $f : \mathcal{X} \to \mathcal{Y}$. La mesure naturelle d'erreur est la perte 0-1 (\emph{0-1 loss}) définie par : 
\begin{equation}
    l(y, \hat{y}) = \mathbbm{1}_{y \neq \hat{y}}.
\end{equation}
La minimisation exacte du risque empirique sous la perte 0-1 sur une classe d'hypothèses générale est un problème combinatoire NP-difficile. De plus, le risque empirique associé est une fonction constante par morceaux de ses paramètres : son gradient est nul presque partout, ce qui rend inapplicables les méthodes d'optimisation par descente de gradient.

Pour lever ce verrou, on reparamétrise le classifieur discret à l'aide d'une fonction à valeurs réelles.

\begin{definition}[Prédicteur par fonction signe randomisée]
    Soit $g : \mathcal{X} \to \R$ une fonction réelle mesurable. On définit le classifieur binaire $f : \mathcal{X} \to \{-1, 1\}$ par $f(x) = \operatorname{sgn}(g(x))$ avec la convention :
    \begin{equation}
        \operatorname{sgn}(u) = 
        \begin{cases}
            1 & \text{si } u > 0, \\
            -1 & \text{si } u < 0, \\
            U \sim \operatorname{Unif}(\{-1, 1\}) & \text{si } u = 0,
        \end{cases}
    \end{equation}
    où $U$ est une variable aléatoire auxiliaire indépendante du couple d'observations $(X, Y) \sim \mathbb{P}$.
\end{definition}

Cette reparamétrisation permet de découpler la \emph{décision brute} (donnée par $\operatorname{sgn}(g(x))$) de la \emph{confiance / distance à la frontière} (quantifiée par l'amplitude $|g(x)|$). On introduit alors la \emph{marge fonctionnelle} $u = Y g(X)$ :
\begin{itemize}
    \item $u > 0$ : la prédiction est correcte ($\operatorname{sgn}(g(X)) = Y$).
    \item $u < 0$ : la prédiction est erronée ($\operatorname{sgn}(g(X)) \neq Y$).
    \item $u = 0$ : le point se situe exactement sur la frontière de décision.
\end{itemize}

\begin{proposition}[Identité de marge pour la perte 0-1]
    Le risque théorique de classification binaire associé à $f = \operatorname{sgn} \circ \, g$ s'exprime comme l'espérance d'une fonction scalaire de la marge fonctionnelle $Y g(X)$ :
    \begin{equation}
        R(f) = \mathbb{E}_{(X,Y) \sim \mathbb{P}} \left[ \Phi_{0-1}(Y g(X)) \right],
    \end{equation}
    avec la fonction de saut :
    \begin{equation}
        \Phi_{0-1}(u) = \mathbbm{1}_{u < 0} + \frac{1}{2} \mathbbm{1}_{u = 0}.
    \end{equation}
\end{proposition}

\begin{proof}
    En explicitant l'espérance conjointe sur $(X,Y)$ et sur la variable de randomisation $U$ :
    \begin{align*}
        R(f) &= \mathbb{E}_{(X,Y), U}\left[ \mathbbm{1}_{f(X) \neq Y} \right] \\
        &= \mathbb{E}\left[ \mathbbm{1}_{g(X) \neq 0} \, \mathbbm{1}_{f(X) \neq Y} \right] + \mathbb{E}\left[ \mathbbm{1}_{g(X) = 0} \, \mathbbm{1}_{f(X) \neq Y} \right] \\
        &= \mathbb{P}(Y g(X) < 0) + \mathbb{E}\left[ \mathbbm{1}_{g(X) = 0} \, \mathbbm{1}_{U \neq Y} \right].
    \end{align*}
    Comme $U$ est indépendante de $(X, Y)$ et uniformément distribuée sur $\{-1, 1\}$, $\mathbb{P}(U \neq Y \mid g(X) = 0) = \frac{1}{2}$. D'où :
    \begin{equation*}
        R(f) = \mathbb{P}(Y g(X) < 0) + \frac{1}{2} \mathbb{P}(g(X) = 0) = \mathbb{E}_{(X,Y) \sim \mathbb{P}} \left[ \Phi_{0-1}(Y g(X)) \right].
    \end{equation*}
\end{proof}

% -------------------------------------------------
\subsection{Risque de substitution et théorème de calibration}

Dans le cas général, lorsque l'on fait face à une fonction de coût $l$ non continue ou non convexe, on définit une fonction convexe $ \Phi$ et une fonction $g$ pour définir un risque de substitution ou \emph{surrogate risk} $R_\Phi(g)$. 

\begin{definition}[Risque de substitution -- Surrogate Risk]
    Soit $\Phi : \R \to [0, \infty[$ une fonction convexe. Le risque de substitution associé à $\Phi$ pour une fonction $g : \mathcal{X} \to \R$ est défini par :
    \begin{equation}
        R_\Phi(g) := \mathbb{E}_{(X,Y) \sim \mathbb{P}} \left[ \Phi(Y g(X)) \right].
    \end{equation}
    Exemples classiques :
    \begin{itemize}
        \item Perte charnière (Hinge) : $\Phi(u) = (1 - u)_+$,
        \item Perte quadratique : $\Phi(u) = (1 - u)^2$,
        \item Perte logistique : $\Phi(u) = \ln(1 + e^{-u})$.
    \end{itemize}
\end{definition}

Or, nous ne sommes pas assurés qu'optimiser sur $R_\Phi(g)$ nous donne une solution optimale $ \hat{f}$ pour notre risque initial $ R(f)$. Le théorème de calibration ci-dessous nous donne une estimation de l'erreur commise. 

\begin{theorem}[Calibration -- Zhang, Bartlett-Jordan-McAuliffe]
    Soit $\Phi : \R \to [0, \infty[$ une fonction convexe, dérivable en 0 avec $\Phi'(0) < 0$. Il existe une fonction convexe, non décroissante et continue $\psi : [0, 1] \to [0, \infty[$ vérifiant $\psi(0) = 0$ telle que pour toute fonction mesurable $g : \mathcal{X} \to \R$ :
    \begin{equation}
        \psi\left( R(\operatorname{sgn} \circ \, g) - R^* \right) \leqslant R_\Phi(g) - \inf_{g'} R_\Phi(g').
    \end{equation}
    En particulier :
    \begin{itemize}
        \item Pour la perte charnière : $\psi(t) = t$,
        \item Pour la perte quadratique : $\psi(t) = t^2$,
        \item Pour la perte logistique : $\psi(t) \geqslant \frac{t^2}{2}$.
    \end{itemize}
\end{theorem}

% ======================================================================
% Décomposition du Risque et Classes Finies

\section{Décomposition du Risque et Classes Finies}

% -------------------------------------------------
\subsection{Décomposition canonique de l'excès de risque}

Dans les chapitres précédents, nous avons décomposé le risque empirique en une erreur \emph{d'approximation} (liée au choix de l'espace de fonction à optimiser $ \mathcal{F}$) et une erreur \emph{d'estimation} (liée à l'échantillon restreint $S$ de taille $n$ que l'on considère, purement statistique) : 
\begin{equation}
    R( \hat{f}_S) - R^* = \underbrace{\left(  R( \hat{f}_S) - \underset{f \in \mathcal{F}}{\inf } R(f) \right)}_{\text{estimation}} + \underbrace{\left( \underset{f \in \mathcal{F}}{\inf} R(f) + R^* \right)}_{\text{approximation}}, 
\end{equation}
avec $ \hat{f}_S$ le minimiseur du risque empirique sur $ \mathcal{F}$. Soit $f^\dagger \in \mathcal{F}$ le meilleur prédicteur théorique possible dans $ \mathcal{F}$. On a donc : 
\begin{equation}
    R( \hat{f}_S) - R^* = R( \hat{f}_S) - R(f^\dagger) + R(f^\dagger) - R^*. 
\end{equation}
Or ici, $ \hat{f}_S$ est calculé sur un échantillon fixé $S$, on ne peut donc pas appliquer la LGN sur la taille de l'échantillon pour obtenir une convergence p.s vers $f^\dagger$. Pour séparer ce qui dépend des données et ce qui dépend de l'estimation dans ce risque, on introduit les risques empiriques $ \hat{R}_S( \hat{f}_S)$ et $ \hat{R}_S(f^\dagger)$ : 
\begin{equation}
    R( \hat{f}_S) - R(f^\dagger) = \left[R( \hat{f}_S) - \hat{R}_S( \hat{f}_S)\right] + \left[ \hat{R}_S( \hat{f}_S) - \hat{R}_S(f^\dagger)\right] + \left[ \hat{R}_S(f^\dagger) - R(f^\dagger)\right]. 
\end{equation}
En regardant chaque élément : 
\begin{itemize}
    \item $ R( \hat{f}_S) - \hat{R}_S( \hat{f}_S)$ correspond à l'erreur commise sur le calcul du risque de l'estimateur restreint à l'échantillon $S$. Comme $ \hat{f}_S \in \mathcal{F}$, on peut le majorer brutalement par le pire estimateur de la classe $ \mathcal{F}$ : 
    \begin{equation}
        R( \hat{f}_S) - \hat{R}_S( \hat{f}_S) \leqslant \underset{f \in \mathcal{F}}{\sup} \left| R(f) - \hat{R}_S(f) \right|.
    \end{equation}
    \item $ \hat{R}_S( \hat{f}_S) - \hat{R}_S(f^\dagger)$ correspond à l'erreur empirique du risque entre l'optimum calculé et le meilleur théorique,
    \item $\hat{R}_S(f^\dagger) - R(f^\dagger)$ correspond à l'erreur commise sur le calcul du risque sur l'échantillon $S$ du meilleur estimateur théorique $f^\dagger$. De la même façon que pour le premier, on peut le majorer par le pire estimateur de $ \mathcal{F}$ : 
    \begin{equation}
        \hat{R}_S( \hat{f}_S) - R(f^\dagger) \leqslant \underset{f \in \mathcal{F}}{\sup} \left| \hat{R}_S(f) - R(f) \right|. 
    \end{equation}
\end{itemize}

Le terme du milieu $ \hat{R}_S( \hat{f}_S) - \hat{R}_S(f^\dagger)$ correspond à l'erreur d'optimisation. Cette erreur est simplement due à la performance de notre algorithme ou à la puissance de notre unité de calcul. On peut donc la majorer arbitrairement par une constante $ \varepsilon_{opt}$ (qui dépend par exemple de la puissance de notre GPU, du temps de calcul, etc\dots). 

\begin{proposition}[Décomposition canonique de l'excès de risque]
    Soit $\mathcal{F}$ une classe d'hypothèses, $\hat{f}_S \in \arg\min_{f \in \mathcal{F}} \hat{R}_S(f)$ le minimiseur empirique et $f^\dagger \in \arg\min_{f \in \mathcal{F}} R(f)$ le meilleur estimateur possible dans $ \mathcal{F}$. L'erreur d'estimation se majore par :
    \begin{equation}
        \boxed{R(\hat{f}_S) - R(f^\dagger) \leqslant 2 \sup_{f \in \mathcal{F}} |\hat{R}_S(f) - R(f)| + \varepsilon_{\mathrm{opt}}}
    \end{equation}
    où $\varepsilon_{\mathrm{opt}} = \hat{R}_S(\hat{f}_S) - \inf_{f \in \mathcal{F}} \hat{R}_S(f)$ représente l'erreur d'optimisation numérique (nulle si la minimisation est exacte).
\end{proposition}

On souhaite maintenant estimer (ou borner) l'erreur maximale que l'on peut faire dans $ \mathcal{F}$ en estimant le risque $ \hat{R}_S$ sur l'échantillon $S$ : $ \sup_{f \in \mathcal{F}} |\hat{R}_S(f) - R(f)|$. On va pour cela distinguer deux cas : si $ \mathcal{F}$ est une classe finie ou non. 

% -------------------------------------------------
\subsection{Majoration de l'erreur du risque -- cas fini}

On suppose ici que $\mathcal{F}$ dispose d'un nombre fini d'éléments $|\mathcal{F}| = M < \infty$ et que les pertes sont bornées par $l_\infty$. Pour majorer cette quantité, on utilise l'inégalité de Hoeffding. 

\begin{theorem}[Inégalité de Hoeffding]
    Soit $(X_k)_{1 \leqslant k \leqslant n}$ une suite de variables aléatoires réelles \emph{indépendantes} telles que pour deux suites $(a_k)_{1 \leqslant k \leqslant n}$ et $ (b_k)_{1 \leqslant k \leqslant n}$, on ait : 
    \begin{equation}
        \forall k \in \llbracket 1, n \rrbracket, \quad \myP(a_k \leqslant X_k \leqslant b_k) = 1. 
    \end{equation}
    Alors, pour tout $ t \geqslant 0$, on peut minorer la concentration de $S_n = X_1 + \dots + X_n$ autour de sa moyenne en probabilité : 
    \begin{equation}
        \myP(|S_n - \mathbb{E}(S_n)| \geqslant t) \leqslant 2 \exp \left( - \frac{2 t^2}{ \sum_{i=1}^{n} (b_i - a_i)^2}\right). 
    \end{equation}
\end{theorem}

\begin{proposition}
    Or ici, pour toute fonction $f \in \mathcal{F}$ fixée indépendante de l'échantillon $S$, la perte sur chaque observation $ Z_i = l(y_i, f(x_i))$ est une variable aléatoire réelle telle que : 
    \begin{enumerate}
        \item Les $(Z_i)_{1 \leqslant i \leqslant n}$ sont i.i.d. 
        \item Pour tout $ i \in \llbracket 1,n \rrbracket$, on a $ Z_i \in [0, l_\infty]$. 
        \item La moyenne empirique est : $ \hat{R}_S(f) = \frac{1}{n} \sum_{i=1}^{n} Z_i$. 
        \item L'espérance est : $ \mathbb{E}[Z_i] = R(f)$. 
    \end{enumerate}
    On peut donc appliquer l'inégalité de Hoeffding sur la moyenne empirique $ \hat{R}_S(f)$ avec $ a_k = 0$ et $ b_k = l_\infty / n$ pour tout $ k \in \llbracket 1, n \rrbracket$. D'où : 
    \begin{equation}
        \forall t \geqslant 0, \quad \myP( |R(f) - \hat{R}_S(f) | \geqslant t) \leqslant 2 \exp \left(- \frac{2 n t^2}{l_\infty ^2}\right). 
    \end{equation}
\end{proposition}

Sachant contrôler chaque estimateur de la classe $ \mathcal{F}$ indépendamment, on souhaite maintenant les contrôler tous\footnote{"Un anneau unique pour les contrôler tous".}.
Pour contrôler simultanément tous les estimateurs de la classe, on applique l'inégalité de Boole (\emph{Union Bound})\footnote{Cette inégalité découle des propriétés des mesures. Elle affirme que pour une famille au plus dénombrable de partie mesurables $A_1, A_2, A_3, \dots$, pour toute mesure $ \mu$, on a : $ \mu \left( \bigcup_n A_n\right) \leqslant \sum_{n} \mu (A_n)$.} sur les $M = |\mathcal{F}|$ éléments :
\begin{align*}
    \myP\left( \sup_{f \in \mathcal{F}} |\hat{R}_S(f) - R(f)| \geqslant t \right) 
    &\leqslant \sum_{f \in \mathcal{F}} \myP\left( |\hat{R}_S(f) - R(f)| \geqslant t \right) \\
    &\leqslant 2 |\mathcal{F}| \exp\left( -\frac{2nt^2}{l_\infty^2} \right).
\end{align*}

En posant cette probabilité d'erreur égale à $\delta \in ]0, 1[$, on obtient le seuil d'écart $t = l_\infty \sqrt{\frac{\ln(2|\mathcal{F}|/\delta)}{2n}}$.

\begin{theorem}[Borne uniforme pour les classes finies]
    Sous les hypothèses $l(y, f(x)) \in [0, l_\infty]$ et $|\mathcal{F}| < \infty$, pour tout $\delta \in ]0, 1[$, avec probabilité au moins $1 - \delta$ sur le tirage de $S \sim \myP^{\otimes n}$ :
    \begin{equation}
        \sup_{f \in \mathcal{F}} |\hat{R}_S(f) - R(f)| \leqslant l_\infty \sqrt{\frac{\ln(2|\mathcal{F}|/\delta)}{2n}}.
    \end{equation}
    Par conséquent, pour le minimiseur empirique exact $\hat{f}_S$, on a avec probabilité au moins $1 - \delta$ :
    \begin{equation}
        R(\hat{f}_S) \leqslant \inf_{f \in \mathcal{F}} R(f) + 2 l_\infty \sqrt{\frac{\ln(2|\mathcal{F}|/\delta)}{2n}} + \varepsilon_{\mathrm{opt}}.
    \end{equation}
\end{theorem}

\begin{remark}
    Le résultat est immédiat ici. Si les hypothèses sont vérifiées, nous sommes capables de borner le risque. En revanche, si $| \mathcal{F}| = \infty$, les bornes de celui-ci explosent. Nous devons donc décrire une nouvelle théorie pour traiter ce cas. 
\end{remark}

% -------------------------------------------------
\subsection{Majoration de l'erreur du risque -- cas infini}

Analysons le résultat précédent : on calcule la probabilité que $ \sup_{f \in \mathcal{F}} |\hat{R}_S(f) - R(f)| \geqslant t $ par une somme finie d'unions d'événements sur $ \mathcal{F}$. C'est donc de là que la borne explose si $ \mathcal{F}$ est infini. 

La solution est donc de traiter la classe entière de fonctions $ \mathcal{F}$ en considérant une application $g$ définie sur l'échantillon $S$ qui lui associe l'erreur la plus importante possible sur la classe de fonctions $ \mathcal{F}$ : 

\begin{equation}
    g(S) := \underset{f \in \mathcal{F}}{\sup} \left( \hat{R}_S(f) - R(f)\right). 
\end{equation}

C'est donc maintenant l'échantillon qui définit la variation d'erreur. On souhaite toujours la borner : on introduit donc la \emph{méthode des différences bornées}. 

\begin{definition}[Différences Bornées]
    Soit $ f : Z_1 \times \dots \times Z_n \longrightarrow \R$ une application. On dit que $f$ vérifie la \emph{condition des différences bornées} pour les constantes $(c_1, \dots, c_n) \in \R^n$ si pour tout $ i \in \llbracket 1, n \rrbracket$ on a : 
    \begin{equation}
        \underset{z_1, \dots, z_i, \dots, z_n, z_i'}{\sup} \left| f(z_1, \dots, z_i, \dots, z_n) - f(z_1, \dots, z_i', \dots, z_n)\right| \leqslant c_i. 
    \end{equation}
    Autrement dit, si l'on change la ième entrée $z_i$ de $f$ en $z_i'$, sa sortie ne bouge pas plus de $c_i$, et cela pour tout $i$. 
\end{definition}

\begin{proposition}[Différences bornées sur $g$]
    Appliquons les différences bornées sur l'application $g$. Puisque pour tout $f \in \mathcal{F}$, le risque $R(f)$ ne dépend pas des données $S$, on a : 
    \begin{equation}
        \hat{R}_S(f) - \hat{R}_{S'}(f) = \frac{1}{n} \left( l(z_i, f) - l(z_i', f)\right). 
    \end{equation}
    Or, par hypothèse, la perte directe est contenue dans $[0, l_\infty]$. Nous devons passer au $\sup$ pour conclure sur $g$. 
\end{proposition}

\begin{lemma}[Passage au $\sup$]
    Soient $u, v : \mathcal{F} \longrightarrow \R$ deux applications. Pour tout $f \in \mathcal{F}$, on a l'implication suivante : 
    \begin{equation}
        \left| u(f) - v(f)\right| \leqslant c \longrightarrow \left| \underset{f \in \mathcal{F}}{\sup} \; u(f) - \underset{f \in \mathcal{F}}{\sup} \; v(f)\right| \leqslant c. 
    \end{equation}
\end{lemma}

En appliquant le lemme précédent avec $c = \frac{l_\infty}{n}$, on en déduit que pour tout $i \in \llbracket 1, n \rrbracket$ : 
\begin{equation}
        |g(S) - g(S')| = \left| \sup_{f \in \mathcal{F}} (\hat{R}_S(f) - R(f)) - \sup_{f \in \mathcal{F}} (\hat{R}_{S'}(f) - R(f)) \right| \leqslant \frac{l_\infty}{n}.
    \end{equation}
    La fonction $g$ vérifie donc la propriété des différences bornées avec les constantes coordonnées :
    \begin{equation}
        c_i = \frac{l_\infty}{n} \quad \text{pour tout } i \in \llbracket 1, n \rrbracket.
    \end{equation}

\begin{theorem}[Inégalité de McDiarmid]
    Soient $Z_1, \dots, Z_n$ des variables aléatoires indépendantes à valeurs dans un espace mesurable $\mathcal{Z}$. Soit $\phi : \mathcal{Z}^n \to \R$ une fonction vérifiant l'hypothèse des différences bornées pour des constantes $(c_1, \dots, c_n) \in \R_+^n$. Alors, pour tout $t \geqslant 0$ :
    \begin{align}
        \myP\left( \phi(S) - \mathbb{E}_S[\phi(S)] \geqslant t \right) &\leqslant \exp\left( -\frac{2t^2}{\sum_{i=1}^n c_i^2} \right), \\
        \myP\left( \mathbb{E}_S[\phi(S)] - \phi(S) \geqslant t \right) &\leqslant \exp\left( -\frac{2t^2}{\sum_{i=1}^n c_i^2} \right), \\
        \myP\left( |\phi(S) - \mathbb{E}_S[\phi(S)]| \geqslant t \right) &\leqslant 2 \exp\left( -\frac{2t^2}{\sum_{i=1}^n c_i^2} \right).
    \end{align}
\end{theorem}

On décompose la déviation uniforme bilatérale en deux déviations unilatérales :
\begin{equation}
    \sup_{f \in \mathcal{F}} |\hat{R}_S(f) - R(f)| \leqslant \sup_{f \in \mathcal{F}} (\hat{R}_S(f) - R(f)) + \sup_{f \in \mathcal{F}} (R(f) - \hat{R}_S(f)).
\end{equation}
Puisque chacune de ces deux applications admet les constantes de variations coordonnées $c_i = \frac{l_\infty}{n}$, la somme des carrés s'évalue explicitement par :
\begin{equation}
    \sum_{i=1}^n c_i^2 = \sum_{i=1}^n \left(\frac{l_\infty}{n}\right)^2 = \frac{l_\infty^2}{n}.
\end{equation}
L'application de la forme unilatérale de McDiarmid fournit alors le contrôle de concentration suivant :

\begin{proposition}[Concentration de la déviation uniforme unilatérale]
    Sous l'hypothèse d'une perte bornée $l(y, f(x)) \in [0, l_\infty]$, pour tout $\delta \in ]0, 1[$, chacune des inégalités suivantes est vérifiée avec probabilité au moins $1 - \delta$ sur le tirage de $S \sim \myP^{\otimes n}$ :
    \begin{align}
        \sup_{f \in \mathcal{F}} (\hat{R}_S(f) - R(f)) &\leqslant \mathbb{E}_S \left[ \sup_{f \in \mathcal{F}} (\hat{R}_S(f) - R(f)) \right] + l_\infty \sqrt{\frac{\ln(1/\delta)}{2n}}, \\
        \sup_{f \in \mathcal{F}} (R(f) - \hat{R}_S(f)) &\leqslant \mathbb{E}_S \left[ \sup_{f \in \mathcal{F}} (R(f) - \hat{R}_S(f)) \right] + l_\infty \sqrt{\frac{\ln(1/\delta)}{2n}}.
    \end{align}
\end{proposition}

\begin{remark}
    Ce résultat marque une avancée théorique majeure : la déviation uniforme se concentre à une vitesse en $\mathcal{O}(1/\sqrt{n})$ autour de son espérance, \emph{sans aucune hypothèse de finitude ou de structure} sur la classe $\mathcal{F}$.
    
    Toutefois, le problème n'est pas encore résolu : la déviation se concentre autour d'une espérance théorique inconnue $\mathbb{E}_S\left[ \sup_{f \in \mathcal{F}} (\hat{R}_S(f) - R(f)) \right]$. Si la classe $\mathcal{F}$ est arbitrairement riche (par exemple l'ensemble de toutes les fonctions mesurables), cette espérance vaut identiquement $1$ pour tout $n$. L'objectif de la section suivante consiste donc à borner cette espérance résiduelle à l'aide de la \emph{symétrisation} et de la \emph{complexité de Rademacher}.
\end{remark}


% ======================================================================
% Complexité de Rademacher

\section{Complexité de Rademacher}

\begin{notation}
    Pour alléger les calculs, on note $z = (x,y)$ et $z_i = (x_i, y_i)$ tirés selon la loi $\myP$. Pour tout prédicteur $f \in \mathcal{F}$ et toute fonction de perte $l : \mathcal{Y} \times \mathcal{Y} \to \R$, on associe la fonction de perte induite $h_f : \mathcal{X} \times \mathcal{Y} \to \R$ définie par $h_f(z) = l(y, f(x))$. On étudie ainsi la classe de fonctions :
    \begin{equation}
        \mathcal{H} = l \circ \mathcal{F} = \{ z \mapsto l(y, f(x)) \mid f \in \mathcal{F} \}. 
    \end{equation}
    Le risque empirique devient $\hat{R}_S(f) = \frac{1}{n} \sum_{i=1}^n h(z_i)$ et le risque théorique $R(f) = \mathbb{E}_z[h(z)]$.
\end{notation}

La section précédente a établi que la déviation uniforme se concentre autour de son espérance théorique :
\begin{equation}
    \mathbb{E}_S \left[ \sup_{f \in \mathcal{F}} (R(f) - \hat{R}_S(f)) \right] = \mathbb{E}_S \left[ \sup_{h \in \mathcal{H}} \left( \mathbb{E}_z[h(z)] - \frac{1}{n} \sum_{i=1}^n h(z_i) \right) \right].
\end{equation}
Le processus de \emph{symétrisation} permet de majorer cette espérance par une quantité calculable à partir de perturbations aléatoires de l'échantillon.

% -------------------------------------------------
\subsection{Le procédé de symétrisation}

\begin{proposition}[Symétrisation par échantillon fantôme]
    Soit $S' = (z_1', \dots, z_n') \sim \myP^{\otimes n}$ un second échantillon (dit \emph{échantillon fantôme}) indépendant de $S$ et de même loi. Par linéarité de l'espérance :
    \begin{equation}
        \mathbb{E}_z[h(z)] = \mathbb{E}_{S'} \left[ \frac{1}{n} \sum_{i=1}^n h(z_i') \right].
    \end{equation}
    Le premier échantillon $S$ étant indépendant de $S'$, on regroupe les termes :
    \begin{equation}
        \mathbb{E}_z[h(z)] - \frac{1}{n} \sum_{i=1}^n h(z_i) = \mathbb{E}_{S'} \left[ \frac{1}{n} \sum_{i=1}^n \big(h(z_i') - h(z_i)\big) \right].
    \end{equation}
    Par convexité de la fonction supremum (inégalité de Jensen, $\sup \mathbb{E} \leqslant \mathbb{E} \sup$), puis en intégrant par rapport à $S$ :
    \begin{equation}
        \mathbb{E}_S \left[ \sup_{h \in \mathcal{H}} \left( \mathbb{E}_z[h(z)] - \frac{1}{n} \sum_{i=1}^n h(z_i) \right) \right] \leqslant \mathbb{E}_{S, S'} \left[ \sup_{h \in \mathcal{H}} \frac{1}{n} \sum_{i=1}^n \big(h(z_i') - h(z_i)\big) \right].
    \end{equation}
\end{proposition}

\begin{proposition}[Injection des variables de Rademacher]
    Posons $d_i = h(z_i') - h(z_i)$. Comme $z_i$ et $z_i'$ sont i.i.d., $d_i \overset{\mathcal{L}}{=} -d_i$. Soient $\varepsilon_1, \dots, \varepsilon_n$ des variables de Rademacher indépendantes de $S, S'$ telles que $\myP(\varepsilon_i = 1) = \myP(\varepsilon_i = -1) = \frac{1}{2}$. Alors $(d_1, \dots, d_n) \overset{\mathcal{L}}{=} (\varepsilon_1 d_1, \dots, \varepsilon_n d_n)$ pour tout $h \in \mathcal{H}$ conjointement :
    \begin{align*}
        \mathbb{E}_{S, S'} \left[ \sup_{h \in \mathcal{H}} \frac{1}{n} \sum_{i=1}^n d_i \right] 
        &= \mathbb{E}_{S, S', \varepsilon} \left[ \sup_{h \in \mathcal{H}} \frac{1}{n} \sum_{i=1}^n \varepsilon_i \big(h(z_i') - h(z_i)\big) \right] \\
        &\leqslant \mathbb{E}_{S', \varepsilon} \left[ \sup_{h \in \mathcal{H}} \frac{1}{n} \sum_{i=1}^n \varepsilon_i h(z_i') \right] + \mathbb{E}_{S, \varepsilon} \left[ \sup_{h \in \mathcal{H}} \frac{1}{n} \sum_{i=1}^n (-\varepsilon_i) h(z_i) \right].
    \end{align*}
    Comme $-\varepsilon_i \overset{\mathcal{L}}{=} \varepsilon_i$, les deux termes sont égaux et l'échantillon fantôme s'élimine :
    \begin{equation}
        \mathbb{E}_S \left[ \sup_{h \in \mathcal{H}} \left( \mathbb{E}_z[h(z)] - \frac{1}{n} \sum_{i=1}^n h(z_i) \right) \right] \leqslant 2 \, \mathbb{E}_{S, \varepsilon} \left[ \sup_{h \in \mathcal{H}} \frac{1}{n} \sum_{i=1}^n \varepsilon_i h(z_i) \right].
    \end{equation}
\end{proposition}